% Part: lambda-calculus
% Chapter: church-rosser
% Section: parallel-beta-reduction

\documentclass[../../../include/open-logic-section]{subfiles}

\begin{document}

\olfileid{lam}{cr}{pb}

\olsection{متوازی $\beta$-تقلیل}

اسیں \emph{متوازی $\beta$-تقلیل} دا خیال پیش کردے آں تے
ثابت کردے آں کہ ایہدے وچ چرچ–روسر خاصیت اے۔

\begin{defn}[متوازی $\beta$-تقلیل، $\bredpar$] \ollabel{defn:bredpar}
  ٹرماں دی متوازی تقلیل ($\bredpar$) استقرا نال ایویں متعین ہوندی اے۔
  \begin{enumerate}
    \item \ollabel{defn:bredpar1} $x \bredpar x$۔
    \item \ollabel{defn:bredpar2} جے $N \xrightarrow{\beta} N'$، تاں $ \lambd[x][N] \bredpar
      \lambd[x][N']$۔
    \item \ollabel{defn:bredpar3} جے $P \bredpar P'$ تے $Q \bredpar Q'$، تاں $PQ \bredpar
      P'Q'$۔
    \item \ollabel{defn:bredpar4} جے $N \bredpar N'$ تے $Q \bredpar Q'$، تاں
      $(\lambd[x][N])Q \bredpar \Subst{N'}{Q'}{x}$۔
  \end{enumerate}
\end{defn}

متوازی $\beta$-تقلیل نال ٹرم دے جِنّے مرضی ریڈیکس اکّو قدم
وچ سکیڑ سکدے آں۔ $\beta$-تقلیل نالوں فرق ایہ اے کہ اک قدم
وچ صرف اوہی ریڈیکس سکیڑے جا سکدے نیں جیہڑے اصل ٹرم وچ
موجود نیں۔ جیہڑے متوازی $\beta$-تقلیل دے نتیجے وچ پیدا ہون،
اوہ اوس قدم وچ نہیں سکیڑے جا سکدے۔ مثلاً ٹرم
$(\lambd[f][fx])(\lambd[y][y])$ دی متوازی $\beta$-تقلیل نال
یا اوہو ٹرم ملدا اے یا $(\lambd[y][y])x$۔ ایس قدم وچ
~$x$ تک نہیں پہنچ سکدے، بھاویں $\beta$-تقلیل نال~$x$ مل
سکدا اے۔ وجہ ایہ اے کہ ایہ ریڈیکس متوازی $\beta$-تقلیل
دے اک قدم پچھوں ای پیدا ہوندا اے۔ متوازی $\beta$-تقلیل
دا دوجا قدم لا کے~$x$ مل جاندا اے۔

\begin{thm}\ollabel{thm:refl}
  $M \bredpar M$۔
\end{thm}

\begin{proof}
  مشق۔
\end{proof}

\begin{prob}
  \olref[lam][cr][pb]{thm:refl} ثابت کرو۔
\end{prob}

\begin{defn}[$\beta$-مکمل تقلیل]\ollabel{defn:bcd}
  \emph{$\beta$-مکمل تقلیل}~$\bcd{M}$، جیہڑی $M$ توں ملدی اے،
  استقرا نال ایویں متعین ہوندی اے۔
  \begin{align}
    \bcd{x} &= x \ollabel{defn:bcd1} \\
    \bcd{(\lambd[x][N])} &= \lambd[x][\bcd{N}] \ollabel{defn:bcd2}\\
    \bcd{(PQ)} &= \bcd{P}\bcd{Q} && \text{جے $P$، $\lambd$ تجرید نہیں}
    \ollabel{defn:bcd3} \\
    \bcd{((\lambd[x][N])Q)} &= \Subst{\bcd{N}}{\bcd{Q}}{x} \ollabel{defn:bcd4}
  \end{align}
\end{defn}

جیویں ناں توں پتا لگدا اے، ٹرم دی $\beta$-مکمل تقلیل
«پوری متوازی تقلیل» اے۔ متوازی $\beta$-تقلیل وچ اسیں
کوئی ریڈیکس اویویں چھڈ سکدے آں۔ مکمل تقلیل وچ اوہناں ساریاں
نوں سکیڑنا لازمی اے۔ ایس لئی $(\lambd[f][fx])(\lambd[y][y])$
دی مکمل تقلیل $(\lambd[y][y])x$ اے۔ اصل ٹرم آپ نتیجہ نہیں اے۔

\begin{editorial}
  ایس تعریف وچ دِقّت ایہ اے کہ اسیں ہلے نہیں دسیا کہ
  ($\lambd$-)ٹرماں اُتے فنکشن بازگشت نال کیویں متعین کردے نیں۔
  ایہ اگّے ٹھیک کراں گے۔
\end{editorial}

\begin{lem}\ollabel{lem:comp}
  جے $M \bredpar M'$ تے $R \bredpar R'$، تاں $\Subst{M}{R}{y}
  \bredpar \Subst{M'}{R'}{y}$۔
\end{lem}

\begin{proof}
  $M \bredpar M'$ دے !!{derivation} اُتے استقرا نال۔
  \begin{enumerate}
    \item آخری قدم \olref{defn:bredpar1} اے۔ مشق۔
    \item آخری قدم \olref{defn:bredpar2} اے۔ فیر $M$ دی شکل
      $\lambd[x][N]$ تے $M'$ دی شکل $\lambd[x][N']$ اے،
      جتھے $N \bredpar N'$۔ اسیں ثابت کرنا اے کہ
      $\Subst{(\lambd[x][N])}{R}{y} \bredpar
      \Subst{(\lambd[x][N'])}{R'}{y}$، یعنی
      $\lambd[x][\Subst{N}{R}{y}] \bredpar
      \lambd[x][\Subst{N'}{R}{y}]$۔ ایہ \olref{defn:bredpar2}
      تے استقرائی مفروضے توں سیدھا ملدا اے۔
    \item آخری قدم \olref{defn:bredpar3} اے۔ مشق۔
    \item آخری قدم \olref{defn:bredpar4} اے۔ $M$ دی شکل
      $(\lambd[x][N])Q$ تے $M'$ دی شکل $\Subst{N'}{Q'}{x}$ اے۔
      اسیں ثابت کرنا اے کہ $\Subst{((\lambd[x][N])Q)}{R}{y} \bredpar
      \Subst{\Subst{N'}{Q'}{x}}{R'}{y}$، یعنی
      $(\lambd[x][\Subst{N}{R}{y}])\Subst{Q}{R}{y} \bredpar
      \Subst{\Subst{N'}{R'}{y}}{\Subst{Q'}{R'}{y}}{x}$۔ ایہ
      \olref{defn:bredpar4} تے استقرائی مفروضے توں ملدا اے۔
  \end{enumerate}
\end{proof}

\begin{prob}
  \olref[lam][cr][pb]{lem:comp} دا ثبوت پورا کرو۔
\end{prob}

\begin{lem}\ollabel{lem:cont}
  جے $M \bredpar M'$، تاں $M' \bredpar \bcd{M}$۔
\end{lem}
\begin{proof}
  $M \bredpar M'$ دے !!{derivation} اُتے استقرا نال۔
  \begin{enumerate}
    \item آخری قاعدہ \olref{defn:bredpar1} اے۔ مشق۔
    \item آخری قاعدہ \olref{defn:bredpar2} اے۔ $M$ دی شکل
      $\lambd[x][N]$ تے $M'$ دی شکل $\lambd[x][N']$ اے، جتھے
      $N \bredpar N'$۔ اسیں دکھاؤنا اے کہ $\lambd[x][N'] \bredpar
      \bcd{(\lambd[x][N])}$، یعنی $\lambd[x][N'] \bredpar
      \lambd[x][\bcd{N}]$، جیویں \olref{defn:bcd2} دسدا اے۔
      ایہ \olref{defn:bredpar2} تے استقرائی مفروضے توں ملدا اے۔
    \item آخری قاعدہ \olref{defn:bredpar3} اے۔ $M$ دی شکل $PQ$
      تے $M'$ دی شکل $P'Q'$ اے، کجھ $P$، $Q$، $P'$ تے $Q'$ لئی
      جیہناں وچ $P \bredpar P'$ تے $Q \bredpar Q'$۔
      استقرائی مفروضے مطابق $P'
      \bredpar \bcd{P}$ تے $Q' \bredpar \bcd{Q}$۔
      \begin{enumerate}
        \item جے $P$ دی شکل $\lambd[x][N]$ اے، کجھ $x$ تے $N$ لئی،
          تاں $P'$ دی شکل $\lambd[x][N']$ ہونی چاہیدی اے، کسے
          $N'$ لئی جیہدے نال $N \bredpar N'$۔ استقرائی مفروضے
          مطابق $N' \bredpar \bcd{N}$ تے $Q' \bredpar \bcd{Q}$۔
          فیر $(\lambd[x][N'])Q' \bredpar
          \Subst{\bcd{N}}{\bcd{Q}}{x}$، \olref{defn:bredpar4} مطابق۔
        \item جے $P$ کوئی $\lambd$ تجرید نہیں، تاں $P'Q' \bredpar
          \bcd{P}\bcd{Q}$، \olref{defn:bredpar3} مطابق۔ سجا
          پاسا $ \bcd{PQ}$ اے، \olref{defn:bcd3} مطابق۔
      \end{enumerate}
    \item آخری قاعدہ \olref{defn:bredpar4} اے۔ $M$ دی شکل
      $(\lambd[x][N])Q$ تے $M'$ دی شکل $\Subst{N'}{Q'}{x}$ اے،
      کجھ $x$، $N$، $Q$، $N'$ تے~$Q'$ لئی جیہناں وچ
      $N \bredpar N'$ تے $Q \bredpar
      Q'$۔ استقرائی مفروضے مطابق $N' \bredpar \bcd{N}$ تے
      $Q' \bredpar \bcd{Q}$۔ \olref{lem:comp} مطابق
      $\Subst{N'}{Q'}{x} \bredpar \Subst{\bcd{N}}{\bcd{Q}}{x}$،
      جیہدا سجا پاسا عین $\bcd{((\lambd[x][N])Q)}$ اے۔
  \end{enumerate}
\end{proof}

\begin{prob}
  \olref[lam][cr][pb]{lem:cont} دا ثبوت پورا کرو۔
\end{prob}

\begin{thm}\ollabel{thm:cr}
  $\bredpar$ وچ چرچ–روسر خاصیت اے۔
\end{thm}

\begin{proof}
  \olref{lem:cont} توں سیدھا ملدا اے۔
\end{proof}

\end{document}
